\section{Related Work}
\label{sec:related}

\para{Agreement within a single model.}
Self-consistency samples several reasoning chains from one model and takes the most common final answer~\citep{wang2023selfconsistency}.
It improves math accuracy, and the vote margin predicts correctness.
Related methods use disagreement between samples to flag hallucinations~\citep{manakul2023selfcheckgpt,kuhn2023semantic}, or use the model's own token probabilities as a confidence score~\citep{kadavath2022know}.
Unlike these works, we measure agreement {\em across} models with different training pipelines, using a single greedy answer per model.
We also compare it directly with each model's own sequence probability as a single-model baseline.

\para{Ensembles and LLM panels.}
Classical ensemble theory says that voting helps when members are accurate and their errors are diverse~\citep{dietterich2000ensemble,kuncheva2003diversity}.
For LLMs, \citet{li2024moreagents} report that accuracy grows with the number of sampled agents, and LLM-Blender ranks and fuses outputs from several models~\citep{jiang2023llmblender}.
\citet{verga2024juries} replace a single large judge with a panel of smaller models from different families and report better agreement with humans.
Building on this line of work, we test whether a plurality vote over a diverse pool beats the {\em best} member, not just the average one, and how the answer changes across domains and ensemble sizes.

\para{Correlated errors across models.}
Recent studies find that LLM errors are strongly correlated, and more so for more capable models or models that share a developer or base model~\citep{goel2025greatmodels,kim2025correlated}.
\citet{goel2025greatmodels} argue that this similarity weakens AI oversight, because a judge model tends to favor models that resemble it.
We study the consequence of correlated errors for consensus as a correctness signal, and we compare error correlation across factual, math and ethics tasks under a format-matched control.

\para{Machine ethics benchmarks.}
\ethics~\citep{hendrycks2021ethics} collects crowd-labeled moral judgements about everyday scenarios. Its ``hard'' test split is adversarially filtered, and some of its items are ambiguous.
Other datasets record several human judgements per item, \eg Scruples~\citep{lourie2021scruples} and MoralChoice~\citep{scherrer2023moralchoice}, which allows a comparison with human disagreement.
We use \ethics because it has a single reference label per item, which lets us define correctness in the same way as for factual tasks.
We then examine where model consensus and the crowd label part ways.

\para{Selective prediction.}
Abstaining on low-confidence inputs trades coverage for accuracy~\citep{geifman2017selective}.
We evaluate inter-model agreement as a label-free abstention rule and report coverage--accuracy curves per task.
